\mode*

\section{Disposition}

\begin{frame}
  \begin{figure}
    \includegraphics[height=0.8\textheight]{fig/book.jpg}
    \caption{Front cover of Necessary Conditions of Learning.}
  \end{figure}
\end{frame}

\begin{frame}
  \begin{enumerate}
    \alert{\item What makes humans human?}
    \item What is to be learned?
    \item Sameness and difference in learning
    \item What does the world look like to others?
    \item The art of learning
    \item Making learning possible
    \item Learning to help others to learn
  \end{enumerate}
\end{frame}

\begin{frame}
  \begin{block}{Ch 1 What makes humans human?}
    \begin{itemize}
      \item What is pedagogy?
      \item How does it relate to learning?
    \end{itemize}
  \end{block}
\end{frame}

\section{Learning as a by-product}

\begin{frame}
  \begin{example}[Learning as a by-product or as an aim]
    \begin{itemize}
      \item Apprenticeship (PhD studies?) is learning as a by-product.
    \end{itemize}
  \end{example}

  \pause

  \begin{example}[\enquote{De-pedagogized} learning]
    \begin{itemize}
      \item Authentic vs unauthentic learning
      \item \enquote{Real life} vs school
    \end{itemize}
  \end{example}
\end{frame}

\begin{frame}
  \begin{example}[The teacher's paradox]
    \blockcquote[p.~13]{NecessaryConditionsOfLearning}{%
      \textins{T}he more the teacher does to enable the students to answer the 
      questions asked, to solve the problems given, the fewer opportunities may 
      be left for the students to learn to understand that which they are 
      expected to learn to understand. The reason is that the more clearly the 
      teacher tells the students what is to be done, the less chance the 
      students get to make the necessary distinctions (for instance, between 
      what is critical and what is not).%
    }
  \end{example}

  \pause

  \begin{remark}[The problem]
    \begin{itemize}
      \item Surface vs deep learning~\cite{DeepSurfaceLearning}.
      \item Study compared US and Chinese teachers in maths.
    \end{itemize}
  \end{remark}
\end{frame}

\begin{frame}
  \begin{onlyenv}<1>
    \begin{remark}
      \blockcquote[p.~21]{NecessaryConditionsOfLearning}{%
        Teachers must somehow organize work in the classroom in order to make 
        learning possible.
        A frequent question posed in educational research is:
        Which is the best arrangement for learning?%
      }
    \end{remark}
  \end{onlyenv}
  \begin{example}[{\cite[p.~21]{NecessaryConditionsOfLearning}}]
    \begin{itemize}
      \item Is learning by yourself better than learning by being taught?
      \item Does homework enhance learning?
      \item Is problem-based learning better than lectures for big classes?
      \item Is individualized learning preferable to group work?
      \item Is project work a good idea?
    \end{itemize}
  \end{example}
  \begin{onlyenv}<2>
    \begin{remark}
      \blockcquote[p.~21]{NecessaryConditionsOfLearning}{%
        The problem with questions of this kind is that they cannot be 
        answered.
        It is not that they cannot be answered \emph{yet}, and it is not 
        because of a scarcity of research funds or a scarcity of good ideas.
        They are simply imponderable because of the degree of generality.%
      }
    \end{remark}
  \end{onlyenv}
\end{frame}

\section{Take away}

\begin{frame}
  \begin{center}
    What is to be learned?\\[1em]
    instead of\\[1em]
    What is to be done?
  \end{center}
\end{frame}

\begin{frame}
  \begin{block}{Ch 2 What is to be learned?}
    \begin{itemize}
      \item What is an intended learning outcome?
      \item What should it be?
      \item It depends on the learner: \emph{critical aspects, critical 
        features} --- the atoms of the learning objective.
    \end{itemize}
  \end{block}
\end{frame}
